\section{上下文感知与群聊参与机制}
\subsection{上下文范围与消息解释}
消息适配保留消息标识、平台、机器人、会话、群、发言者、时间、提及与引用关系。Web 预览从服务端获取同账号、同群的近期历史；主动参与执行器读取当前群窗口。Runtime 继续按预算整理这些资料，区分本轮问题、被引用内容和历史发言。

本轮指令决定要完成的任务，历史发言解释指代和讨论背景。例如，旧消息提出周五开会，后续更正为周六，本轮追问“最终时间”，应依据有效更正形成答案。来源于其他群的引用文本可以作为本轮提供的资料，但不能成为读取其他群完整历史的授权。

\subsection{规则感知与模型感知合并}
规则适合识别明确提及、已知业务词和输入格式；模型适合处理课程别称、口语、省略与开放式需求。\code{HybridPerceptionEngine} 将两者组织为结构化感知结果，模型输出经过本地 Schema 校验后才能参与决策。当前感知并非向后续层直接透传任意自然语言推断。

结构化结果用于表达意图、实体、引用、偏好能力类别及置信信息。明确的“评课社区”提示来源，课程实体提示查询对象；“最多三项”属于结果数量约束。模型只提供语义候选，最终可用能力、执行预算和群权限仍由程序检查。

复杂度评估与社交决策分别使用感知信号。当前组合采用确定性策略，不为每一个阶段额外调用一个大模型。高复杂度档位选择需要相应信号和置信条件，低置信、指代不明或任务缺项可以进入澄清流程。

\subsection{Social Engine 的动作选择}
Social Engine 将对内容的理解转化为参与动作。当前核心策略包括直接回应、调用业务能力、澄清和保持静默等行为路径；输入权限、群模式与当前请求条件影响选择。它与外围自动参与执行器职责不同：前者决定本轮怎么处理，后者先判断普通群消息是否值得启动这一轮。

显式提及和自动参与使用不同入口。用户明确提出任务时，系统进入请求处理；自动参与则在目标群已开启的前提下检查事件类型、概率、冷却、每小时上限及是否有进行中的运行。随机抽样只影响参与机会，不应改变事实或服务权限。

\begin{table}[htbp]\centering
\caption{自动参与的上下文窗口配置（当前代码值）}
\begin{tabularx}{\textwidth}{@{}L{32mm}L{28mm}L{30mm}Y@{}}\toprule
窗口 & 最大消息数 & 字节预算 & 适用含义\\\midrule
compact & 12 & 3,500 & 紧凑讨论片段\\
standard & 30 & 6,500 & 一般群聊上下文\\
extended & 100 & 7,500 & 更多发言记录，仍受总预算约束\\\bottomrule
\end{tabularx}
\end{table}

窗口越大不等于可完整读入更多长消息，也不等于全天历史。当前控制器维护发送时间和运行占用状态；冷却可以抑制短时连续触发，但不能独立证明跨冷却窗口的建议不会语义重复。这一问题需要专门的长期群聊评估。

\subsection{受授权的能力自适应启用}
当前自适应选择器是可解释的规则实现。它只在管理员允许的只读插件集合中匹配需求；并要求足够的连续讨论、多位发言者和当前明确请求。核对时的条件包括：
\begin{enumerate}
\item 历史至少 8 条，解析得到至少 3 个发言者标签；
\item 最后一条消息包含查询、建议、推荐或疑问等请求信号；
\item 最近 20 条中至少 3 条与同一候选主题相关，最近 4 条仍出现该主题；
\item 候选插件已经获得本群自适应启用授权。
\end{enumerate}

规则满足后，策略存储记录能力启用和原因，当前请求进入正式 Runtime，控制台读取实际生效状态。该机制能够在讨论逐渐形成需求时开放已有能力；它不下载新插件，不修改程序，也不是在线训练或自主生成 Skill。

\begin{figure}[H]\centering
\begin{tikzpicture}[node distance=6mm]
\node[box,text width=13cm] (a) {群授权 + 连续讨论窗口 + 当前请求};
\node[box,text width=13cm,below=of a] (b) {规则匹配：发言者标签、主题密度、近期相关性、允许的只读插件};
\node[box,text width=13cm,below=of b] (c) {记录启用与原因 $\rightarrow$ Runtime 查询 $\rightarrow$ Web 状态同步};
\draw[flow](a)--(b);\draw[flow](b)--(c);
\end{tikzpicture}
\caption{当前自适应启用机制。启用范围来自管理员配置，语义回答由后续 Runtime 完成。}
\end{figure}

四人 20 条固定讨论证明了这一链路的代表场景：前 19 条评课关闭，第 20 条形成请求后启用评课。20 条是该场景的设置，不是算法对所有群的固定阈值；更多群聊类型、话题切换和误触发仍需单独评估。
\source{[S3, S4, S10]；\nolinkurl{adapters/proactive_talk.py}、\nolinkurl{adapters/adaptive_plugins.py} 与 \nolinkurl{perception/social.py}。}
